\SourceLeaf{362}
\Article{7}
Die erste Methode gibt nun folgendes Resultat:
\[
\begin{aligned}
4\text{ Decident}={}&\mathrm{I}\quad\Sigma\varepsilon\text{ von allen}\\
&\phantom{\mathrm{I}\quad}{-}4\Sigma\varepsilon\text{ von denen, wo }y\text{ ganz }[x]\text{ gerade}\\[3pt]
&\mathrm{II.}\ {-}\Sigma\varepsilon\text{ von allen}\\
&\phantom{\mathrm{II.}\ }{+}4\Sigma\varepsilon\text{ von denen, wo }[x]\text{ gerade }[y]\text{ ungerade}\\[3pt]
&\mathrm{IV.}\ {+}4\Sigma\varepsilon\text{ von denen, wo nicht zugleich}\\
&\phantom{\mathrm{IV.}\ {+}4\Sigma\varepsilon\text{ von denen, wo }}[x]\text{ und }[y]\text{ gerade}\\
&\phantom{\mathrm{IV.}\ }{+}Q
\end{aligned}
\]

Hier ist
\[
\begin{array}{cc|l}
\text{für }A&B&Q=\\
+&+&-\text{Intens.}\ {+}\mu\omega i+\text{Int.}\ \frac12m-\mu\omega\\
+&-&-\text{Intens.}\ {+}\mu\omega+\text{Int.}\ \frac12mi-\mu\omega i\\
-&-&-\text{Intens.}\ {-}\mu\omega i-\text{Int.}\ \frac12m+\mu\omega\\
-&+&-\text{Intens.}\ {-}\mu\omega-\text{Int.}\ \frac12mi+\mu\omega i
\end{array}
\]
folgende Tabelle stellt dies dar

\[
\setlength{\arraycolsep}{4pt}
\begin{array}{cc|rr|rr|rr|rr}
\alpha&\beta& A&B&A&B&A&B&A&B\\
&&+&+&+&-&-&-&-&+\\ \cline{3-10}
+&+&+2&0&0&+2&-3&-5&-3&-5\\
+&-&0&+2&-2&0&-5&-3&-1&-3\\
-&-&-3&-1&-1&+1&-4&-2&0&-2\\
-&+&+1&-1&-1&+1&0&-2&-4&-6\\[3pt]
\multicolumn{2}{r}{\dfrac{m-1}{2}}&\text{par,}&\text{impar}&\multicolumn{6}{c}{}
\end{array}
\]

Pars prior ipsius
\[
2Q\dotstwo -3-(\alpha\beta AB)+(A\alpha)+(A\beta)+(B\alpha)-(B\beta)
\]
Das ganze $2Q$ für
\[
\begin{aligned}
\frac{m-1}{2}\text{ impar}\quad&-3+4(A)-(B)-(\beta)+(\alpha A)+(\alpha B)\\
&\quad +(\beta A)-(\beta B)-(\alpha\beta AB)\\[3pt]
\frac{m-1}{2}\text{ par}\quad&-3+2(A)+(B)+(\beta)+(\alpha A)+(\alpha B)\\
&\quad +(\beta A)-(\beta B)+2(\beta AB)-(\alpha\beta AB)
\end{aligned}
\]

\SourceLeaf{363}
Beispiel
\[
\begin{array}{r|r}
m=+3-2i&13\\
M=-1-6i&37
\end{array}
\qquad
\begin{gathered}
(1+i)M+(-2+i)m=1\\
mM=-15-16i
\end{gathered}
\]
\[
\mu=+9+20i
\]

\[
\begin{array}{r|r|r|r}
+4+13i&+38+31i&+1&-1\\
+5+7i&+29+11i&0&+0\\
+6+i&+20-9i&0-i&-3\\
+10+14i&+58+22i&+1&+1\\
+11+8i&+49+2i&+1&-1\\
+12+2i&+40-18i&+1-i&+2\\
+16+15i&+78+13i&+2&-0\\
+17+9i&+69-7i&+1-i&+2\\
+18+3i&+60-27i&+1-i&-2\\ \cline{3-4}
\multicolumn{2}{c}{}&\multicolumn{2}{c}{\text{Decident}=-2}
\end{array}
\]
\[
\begin{array}{lrr}
\mathrm{I}\dotstwo&+1&0\\
\mathrm{II}\dotstwo&0&-4\\
\mathrm{IV}\dotstwo&&0\\
Q=&-5&\\ \hline
\multicolumn{3}{c}{-8}\\
\multicolumn{3}{c}{\text{Gut.}}
\end{array}
\]

\noindent I
\[
\begin{array}{r|r|r|r|r|r|r}
x&X&20\xi&20\eta&20y&20Y&\varepsilon\\
+1&0&6&1&-9&-37&-1\\
&+1&4&4&+4&-28&-1\\
&+2&2&7&+17&-19&-1\\
+2&+1&10&5&-5&-65&-1\\
&+2&8&8&+8&-56&+1\\
\multicolumn{6}{c}{}&-3
\end{array}
\]
\[
\begin{array}{r|r|r|r|r|r|r}
y&Y&20\xi&20\eta&20x&20X&\varepsilon\\
0&-1&3&2&+13&+9&+1^{*}\\
&-2&6&4&+26&+18&+1\\
&-3&9&6&+39&+27&+1\\
+1&-1&2&8&+22&+46&+1\\
\multicolumn{6}{c}{}&+4\\
\multicolumn{6}{c}{}&(+1)
\end{array}
\]

\[
\begin{array}{r|r|r|r|r|r|r}
y&X&9\xi&9\eta&9x&9Y&\varepsilon\\
0&+1&3&2&+13&-20&+1\\
-2&+2&0&3&+6&-3&-1^{*}\\
\multicolumn{6}{c}{}&0(-1)
\end{array}
\]
\[
\begin{array}{r|r|r|r|r|r|r}
x&Y&9\xi&9\eta&9y&9X&\varepsilon\\
+1&-1&1&3&+7&+17&-1\\
+2&-3&4&3&+1&+14&+1\\
\multicolumn{6}{c}{}&0
\end{array}
\]
\[
\mathrm{I}\pdot\pdot\pdot\;+1(0)
\]

\noindent II
\[
\begin{array}{r|r|l|r|r|r}
Y&6\xi&6x&6y&6X&\varepsilon\\
-1&+1&+3+9\omega&-2&-1&-1^{*}\\
-2&+2&+6+9\omega&-4&-2&+1\\
\multicolumn{5}{c}{}&0(-1)
\end{array}
\]
\noindent III
\[
\begin{array}{r|r|r|r|l|r}
x&3\xi&3y&3X&3Y&\varepsilon\\
+1&1&-2&-1&-6+9\omega&-1
\end{array}
\]

\noindent IV
\[
\begin{array}{r|r|l|r|r|r}
X&12\eta&12x&12y&12Y&\varepsilon\\
0&+1&+20+20\omega&-9&-37&-1^{*}\\
+1&+3&+24+20\omega&-3&-39&+1^{*}\\
+2&+5&+28+20\omega&+3&-41&-1\\
\multicolumn{5}{c}{}&-1(0)
\end{array}
\]
\noindent V
\[
\begin{array}{r|r|r|l|r|r}
x&4\eta&4y&4X&4Y&\varepsilon\\
+2&+1&-1&+4-20\omega&-13&+1
\end{array}
\]
\[
\begin{array}{r|r|r|r|r|r}
y&6\eta&6x&6X&6Y&\varepsilon\\
0&+2&+13&+9&-20&+1
\end{array}
\]

\SourceLeaf{364}
\Article{8}
Wir wollen nunmehro das Resultat von I näher betrachten. Es sind vier Combinationen\par
\noindent $1^0$, wenn $x$ und $X$ Ganze sind. Man hat hier
\[
\begin{aligned}
\beta\xi&=-Bx+bX\\
\beta\eta&=-Ax+aX\\
\beta y&=-\alpha x+dX\\
\beta Y&=-Dx+\alpha X
\end{aligned}
\]

Es seien $y^0$ und $Y^0$ die absolut kleinsten Reste von $-\alpha x+dX$, $-Dx+\alpha X$ nach dem Modulus $\beta$ und $\beta y=\beta y'+y^0$, $\beta Y=\beta Y'+Y^0$ und man setze
\[
\begin{alignedat}{2}
\beta u&=-By^0+bY^0\qquad&y^0&=+au-bt\\
\beta t&=-Ay^0+aY^0&Y^0&=+Au-Bt
\end{alignedat}
\]
so werden $t,u$ ganze Zahlen sein, nemlich
\[
\begin{aligned}
-u+ti&=M(x+y'i)-m(X+Y'i)\;{}^*\\
i(t+ui)&=Mi(y'-y)-mi(Y'-Y)
\end{aligned}
\]

und man hat dann $\varepsilon=-1$, wenn\par
\noindent $t+u$ gerade, $\beta,y^0,Y^0$ positiv, oder wenn zwei oder keine Bedingung gilt, sonst $\varepsilon=+1$\par
Wir setzen
\[
\begin{array}{rl}
t+ui=+\theta&\text{wenn }y^0\text{ und }Y^0\text{ beide positiv}\\
-\theta&y^0\text{ und }Y^0\text{ beide negativ}\\
+\theta'&y^0\text{ positiv }Y^0\text{ negativ}\\
-\theta'&y^0\text{ negativ }Y^0\text{ positiv}
\end{array}
\]
\[
\begin{array}{llr}
\text{jedem durch }1+i\text{ theilbaren}&\theta&\text{entspricht dann ein }\varepsilon=-1\\
\text{jedem durch }1+i\text{ theilbaren}&\theta'&\varepsilon=+1\\
\text{jedem durch }1+i\text{ untheilbaren}&\theta&\varepsilon=+1\\
\text{jedem durch }1+i\text{ untheilbaren}&\theta'&\varepsilon=-1
\end{array}
\]
insofern $\beta$ positiv.\par
Für negative $\beta$ ist es umgekehrt.

\SourceLeaf{365}
\[
\setlength{\arraycolsep}{4pt}
\begin{array}{r|r|r|r|r|r|r|r}
x&X&y^0&Y^0&t+ui&\theta&\theta'&\varepsilon\\
+1&0&-9&+3&0-3i&&0+3i&-1\\
&+1&+4&-8&-1+2i&&-1+2i&-1\\
&+2&-3&+1&0-i&&0+i&-1\\
+2&+1&-5&-5&-1-i&+1+i&&-1\\
&+2&+8&+4&+1+2i&+1+2i&&+1
\end{array}
\]

\Article{9}
$2^0$, wenn $y$ und $Y$ Ganze. Es seien hier $x',X'$ die nächsten Ganzen bei $x$ und $X$, und
\[
x-x'=\frac{x^0}{\beta},\quad X-X'=\frac{X^0}{\beta}
\]
und man setze
\[
\begin{gathered}
\beta(t+ui)=-Mx^0+mX^0=-Mp^0+mP^0,\\
t+ui=Mp'-mP'
\end{gathered}
\]
d. i.
\[
\begin{alignedat}{2}
\beta t&=-Ax^0+aX^0\qquad&x^0&=-bt+au\\
\beta u&=-Bx^0+bX^0&X^0&=-Bt+Au
\end{alignedat}
\]

Man hat dann
\[
\begin{aligned}
\varepsilon=-1,&\text{ wenn }\beta\text{ positiv, }x^0\text{ positiv, }X^0\text{ positiv, }t+u\text{ gerade}\\
&\text{etc.}
\end{aligned}
\]
Wir setzen
\[
\begin{array}{rll}
&&\textit{besser}\\
t+ui=+\theta&\text{wenn }x^0\text{ und }X^0\text{ positiv}&+\theta''\\
=-\theta&\text{wenn beide negativ}&-\theta''\\
=+\theta'&\text{wenn }x^0\text{ positiv, }X^0\text{ negativ}&-\theta\\
=-\theta'&\text{wenn }x^0\text{ negativ, }X^0\text{ positiv}&+\theta
\end{array}
\]
wo für $\varepsilon$ dieselbe Regel gelten wird wie oben

\[
\setlength{\arraycolsep}{4pt}
\begin{array}{r|r|r|r|r|r|r|r}
y&Y&x^0&X^0&t+ui&\theta&\theta'&\varepsilon\\
0&-1&-7&+9&+1-3i&&-1+3i&+1\\
&-2&+6&-2&0+2i&&0+2i&+1\\
&-3&-1&+7&+1-i&&-1+i&+1\\
+1&-1&+2&+6&+1&+1&&+1
\end{array}
\]

\SourceLeaf{366}
Man kann nun beweisen
\begin{enumerate}\renewcommand{\labelenumi}{\arabic{enumi})}
\item Dass alle $\theta$, die aus (1) und aus (2) hervorgegangen sind, unter einander verschieden sind. Ihr Complexus heisse $\Theta$.
\item Dass alle $\theta=T+Ui$ die Eigenschaft haben, dass
\[
\begin{gathered}
-bT+aU\\
-BT+AU
\end{gathered}
\]
positive Zahlen kleiner als $\frac12\beta$ sind
\item Dass wenn $T,U$ zwei der eben genannten Bedingungen unterworfene ganze Zahlen sind, $T+Ui$ sich gewiss in $\Theta$ findet.\par
(wie denn? es wird auf obige Gleichung $^*$ gegründet.)
\item Auf ähnliche Weise verhält es sich mit $\theta'$, deren Complexus $\Theta'$ aus denjenigen Zahlen $T'+U'i$ bestehen wird, für welche
\[
\begin{gathered}
-bT'+aU'\\
-(-BT'+AU')
\end{gathered}
\]
positive Zahlen kleiner als $\frac12\beta$.
\end{enumerate}

In unserm Falle ist
\[
\setlength{\arraycolsep}{4pt}
\begin{array}{r|r||r|r||r|r||r|r}
\theta&\varepsilon&\theta'&\varepsilon&\theta''&\varepsilon&\theta'''&\varepsilon\\
+1&+1&0+i&-1&+2-i&+1&+1&-1\\
+1+i&-1&0+2i&+1&+3-i&-1&+2&+1\\
+1+2i&+1&0+3i&-1&\multicolumn{4}{c}{}\\
\multicolumn{2}{c||}{}&-1+i&+1&\multicolumn{4}{c}{}\\
\multicolumn{2}{c||}{}&-1+2i&-1&\multicolumn{4}{c}{}\\
\multicolumn{2}{c||}{}&-1+3i&+1&\multicolumn{4}{c}{}
\end{array}
\]

Hier ist
\[
\begin{aligned}
\theta&=-{}_*M+{}_*m\\
\theta'&=-{}_*M-{}_*m\\
\theta''&=+{}_*Mi+{}_*m\\
\theta'''&=+{}_*Mi-{}_*m
\end{aligned}
\]

\SourceLeaf{367}
\Article{10}
$3^0$. $y$ und $X$ Ganze. Es seien $x',Y'$ die nächsten Ganzen bei $x$ und $Y$, und
\[
x'+yi=p',\quad X+Y'i=P';\quad p-p'=\frac{p^0}{\alpha},\quad P-P'=\frac{P^0}{\alpha}
\]
und man setze
\[
\begin{aligned}
i(t+ui)&=Mp'-mP'=-\frac{Mp^0}{\alpha}+\frac{mP^0}{\alpha}\\
&=-\frac{Mx^0}{\alpha}+\frac{miY^0}{\alpha}\quad\text{d. i.}
\end{aligned}
\]
\[
\begin{alignedat}{3}
\alpha t&=-Bx^0+aY^0\quad&\text{so ist}\quad&x^0&=-bt+au\\
\alpha u&=+Ax^0+bY^0&&Y^0&=+At+Bu
\end{alignedat}
\]

Man hat dann
\[
\varepsilon=-1,\text{ wenn }\alpha\text{ positiv, }x^0\text{ positiv, }Y^0\text{ positiv, }t+u\text{ gerade etc.}
\]
Wir setzen
\[
\begin{array}{rll}
&&\textit{besser}\\
t+ui=+\theta''&\text{wenn }x^0\text{ positiv, }Y^0\text{ positiv}&+\theta'\\
=+\theta'''&\text{wenn }x^0\text{ positiv, }Y^0\text{ negativ}&-\theta'''\\
-\theta''&\text{wenn }x^0\text{ negativ, }Y^0\text{ negativ}&-\theta'\\
-\theta'''&\text{wenn }x^0\text{ negativ, }Y^0\text{ positiv}&+\theta'''
\end{array}
\]
\[
\begin{array}{lr}
\text{Es wird also für jedes}&\theta''\pdot\pdot\pdot\;\varepsilon=-1\\
&\theta'''\pdot\pdot\pdot\;\varepsilon=+1
\end{array}
\]
insofern $\theta''$ oder $\theta'''$ durch $1+i$ theilbar und $\alpha$ positiv.

\[
\setlength{\arraycolsep}{4pt}
\begin{array}{r|r|r|r|r|r|r|r}
y&X&x^0&X^0&t+ui&\theta''&\theta'''&\varepsilon\\
0&+1&+4&-2&+2&&+2&+1\\
+1&+2&-3&-3&-3+i&+3-i&&-1
\end{array}
\]

\Article{11}
$4^{\mathrm{te}}$ Classe $x$ und $Y$ Ganze. Nach ähnlichen Praemissen wie in 3 setze man
\[
\begin{aligned}
-(t+ui)&=Mp'-mP'=-\frac{Mp^0}{\alpha}+\frac{mP^0}{\alpha}\\
&=-\frac{Miy^0}{\alpha}+\frac{mX^0}{\alpha}
\end{aligned}
\]
\[
\begin{alignedat}{2}
\alpha t&=-By^0-aX^0\qquad&y^0&=-bt+au\\
\alpha u&=+Ay^0-bX^0&X^0&=-At-Bu
\end{alignedat}
\]
Man hat dann
\[
\varepsilon=+1\text{ wenn }\alpha\text{ positiv, }y^0\text{ positiv, }X^0\text{ positiv, }t+u\text{ gerade etc.}
\]

\SourceLeaf{368}
Wir setzen
\[
\begin{array}{rl}
t+ui=+\theta''&\text{wenn }y^0\text{ positiv }X^0\text{ negativ}\\
+\theta'''&\text{wenn }y^0\text{ positiv }X^0\text{ positiv}\\
-\theta''&\text{wenn }y^0\text{ negativ }X^0\text{ positiv}\\
-\theta'''&\text{wenn }y^0\text{ negativ }X^0\text{ negativ}
\end{array}
\]
für $\varepsilon$ gilt dann die Regel, dass (wie oben in 3), insofern $\alpha$ positiv
\[
\begin{aligned}
\varepsilon=-1\text{ für jedes durch }1+i&\left\{\begin{array}{ll}\text{theilbare}&\theta''\\\text{untheilbare}&\theta'''\end{array}\right.\\[3pt]
\varepsilon=+1\text{ für jedes durch }1-i&\left\{\begin{array}{ll}\text{theilbare}&\theta'''\\\text{untheilbare}&\theta''\end{array}\right.
\end{aligned}
\]

\[
\setlength{\arraycolsep}{4pt}
\begin{array}{r|r|r|r|r|r|r|r}
x&Y&y^0&X^0&t+ui&\theta''&\theta'''&\varepsilon\\
+1&-1&-2&-1&-1&&+1&-1\\
+2&-3&+1&-4&+2-i&+2-i&&+1
\end{array}
\]

Der Complexus aller $\theta''$ aus 3 und 4, den wir durch $\Theta''$ bezeichnen, besteht also aus allen Zahlen $T+Ui$, wofür
\[
\left.\begin{array}{r}
-bT+aU\\
\text{und }+AT+BU
\end{array}\right\}\text{ positiv und kleiner als }\frac12\alpha
\]
der Complexus aller $\theta'''\pdot\pdot\pdot\;(\Theta''')$ aus denen, wo
\[
\left.\begin{array}{r}
-bT+aU\\
-AT-BU
\end{array}\right\}\text{ positiv und kleiner als }\frac12\alpha
\]

\Article{12}
Nach obiger Verbesserung heisst also die Regel so:\par
Es enthalte
\[
\begin{aligned}
\Theta\quad&\text{alle Zahlen }T+Ui,\text{ wo }+bT-aU\text{ positiv}\\
&\quad -BT+AU\text{ positiv und }<\tfrac12\beta\\
\Theta'\quad&\text{alle Zahlen }T+Ui,\text{ wo }-bT+aU\text{ positiv}\\
&\quad +AT+BU\text{ positiv und }<\tfrac12\alpha\\
\Theta''\quad&\text{alle Zahlen }T+Ui,\text{ wo }-bT+aU\text{ positiv}\\
&\quad -BT+AU\text{ positiv und }<\tfrac12\beta\\
\Theta'''\quad&\text{alle Zahlen }T+Ui,\text{ wo }+bT-aU\text{ positiv}\\
&\quad +AT+BU\text{ positiv und }<\tfrac12\alpha
\end{aligned}
\]
insofern resp. $\beta$ oder $\alpha$ positiv.

\SourceLeaf{369}
\[
\begin{array}{c|c|l}
\text{Für alle durch}&\text{ist dann}&\text{wenn}\\
1+i\text{ theilbaren}&\varepsilon=&\\
\theta&+1&\beta\text{ positiv}\\
\theta'&-1&\alpha\text{ positiv}\\
\theta''&-1&\beta\text{ positiv}\\
\theta'''&+1&\alpha\text{ positiv}
\end{array}
\]

\[
\setlength{\arraycolsep}{4pt}
\begin{array}{r|r||r|r||r|r||r|r}
\multicolumn{2}{c||}{\Theta}&\multicolumn{2}{c||}{\Theta'}&\multicolumn{2}{c||}{\Theta''}&\multicolumn{2}{c}{\Theta'''}\\
0-i&-1&+2-i&+1&+1&+1&-1&-1\\
0-2i&+1&+3-i&-1&+1+i&-1&-2&+1\\
0-3i&-1&\multicolumn{2}{c||}{}&+1+2i&+1&\multicolumn{2}{c}{}\\
+1-i&+1&\multicolumn{6}{c}{}\\
+1-2i&-1&\multicolumn{6}{c}{}\\
+1-3i&+1&\multicolumn{6}{c}{}
\end{array}
\]

\begin{center}
\resizebox{.86\textwidth}{!}{\input{diagrams/native369_lattice.tex}}
\end{center}
